\scanpage{57}
\mspage{63}

On en déduit que \deleted{\unsure{la}} \inserted{$f:X\longrightarrow Y$ est une} $\mathrm{HOT}$-équiv. (au sens de la def. 1) ssi c'est une \inserted{hot-}équivalence (ou ``équivalence faible'') i.e. ssi

\[ \mathrm{hot}(f):\mathrm{hot}(X)\xrightarrow{\sim}\mathrm{hot}(Y) \]

est un iso.

\deleted{\unsure{B)}} \uline{Proposition 2} Soit $f:X^{\deleted{\unsure{\prime}}}\xrightarrow{f}\deleted{\unsure{X}}Y$. Conditions équivalentes (si $\unsure{\mathbb D|\underline D}$ dérivateur \uline{à gauche})

a) $f$ est une $\mathbb D$-équiv., et le reste \deleted{pour} \inserted{après} tt chgt. de base $Y'\longrightarrow Y$ qui est un \deleted{\unsure{immersion}} \inserted{isom. local.} \inserted{(p.ex. pour $Y_{/F}\longrightarrow Y$, $F\in\mathrm{Ob}\,Y^{\unsure{\wedge|\prime}}$}

% slanted
\margin{Autre condition équivalente d) Pour tt faisceau \uline{constant} $M_Y$ sur $Y$, on a $M_Y\xrightarrow{\sim}f_*f^*M_Y$ d$'$) Comme d), avec faisceau \uline{loc} constant}

b) \deleted{$f$ \unsure{est une} $\mathbb D$-\unsure{équiv.} $f$} Pour tout $y\in\mathrm{Ob}\,Y$, le foncteur induit $X_{/y}\longrightarrow Y_{/y}$ est une $\mathbb D$-équiv.

c) Pour tout $y\in\mathrm{Ob}\,Y$, $X_{/y}$ est $\mathbb D$-contractile.

\deleted{\unsure{Défin}} N.B $\rightarrow$ On dit que $Z\xrightarrow{g}Y$ est un \uline{iso. local} si $\forall z\in\mathrm{Ob}\,Z$, posant $y=g(z)$, le morphisme induit pour les \deleted{\unsure{cat}} catégories localisées en $z$ et $y$

\[ Z_{/z}\xrightarrow{\sim}Y_{/y} \]

est un iso. Par exemple toute immersion ouverte, tout \deleted{\unsure{morphisme}} \inserted{foncteur} induit $Y_{/y}\longrightarrow Y$ (foncteur de localisation) \deleted{\unsure{sont}} des iso. loc.

Il en résulte a) $\Rightarrow$ b) trivial$^{\mathrm{mt}}$.

b) $\Leftrightarrow$ c). Résulte du cor. de la page précédente, si on sait que \fbox{$Y_{/y}$ est $\mathbb D$-contractile}. Cela \unsure{résultera|résulte} \deleted{de} Der 5 ci-dessous

\margin{$\searrow$ Cf cor. 1 de Der. 5 (g)}
% archive sheet number: 56
